\subsection*{同步与异步 \texttt{parallel-by}}

和 \co@inthreads@ 类似，\co@parallel-by@ 也可以加上 \co@.async@ 变成异步：

\begin{lstlisting}[language=croqtile,caption={一次异步启动}]
parallel.async p by 6 : block {
}
\end{lstlisting}

不加 \co@.async@ 时，主机会等核函数结束后才继续：

\begin{lstlisting}[language=croqtile,caption={同步形式}]
parallel p by 6 : block {
}
// host continues only after the kernel finishes
\end{lstlisting}

加上 \co@.async@ 之后，主机代码在启动核函数后立刻继续。要让计算和数据传输重叠，或者并发启动彼此独立的核函数，这一点很关键。
